\section[Conclusão]{Conclusão}

O projeto ENSportive foi uma experiência extremamente enriquecedora. Ingressei na
\ac{ages} I sem experiência prévia em desenvolvimento de software e, ao longo do
semestre, precisei aprender tecnologias novas em ritmo acelerado e colocá-las em
prática em um ambiente profissional real. Esse processo me fez crescer muito tanto
tecnicamente quanto no aspecto comportamental.

\textbf{Aprendizados técnicos e de soft skills:}
O ambiente da \ac{ages} é extremamente propício para o desenvolvimento de
\textit{soft skills}. A quantidade de conhecimento novo, as tecnologias desconhecidas
e os prazos curtos exigiram adaptação constante, comunicação eficaz e colaboração
genuína com o time. Aprendi a pedir ajuda no momento certo, a aceitar críticas
construtivas e a contribuir com colegas mais experientes de forma proativa.

\textbf{Sobre o projeto:}
A ENSportive foi um desafio do início ao fim, dada a série de regras de negócio
específicas que precisavam ser respeitadas. O projeto foi construído de forma muito
profissional, contando com desenvolvedores plenos no time, o que elevou o padrão
das entregas. Por mais desgastante que o semestre tenha sido, por mais que tenhamos
precisado fazer cortes no escopo, conseguimos entregar e realizar uma boa
apresentação final. Acredito que nossa stakeholder tenha ficado satisfeita, pois o
sistema poderá ajudar de forma significativa a Escola de Raquetes ENSportive.

\textbf{Reflexão sobre a competição:}
Gostaria de registrar uma reflexão sobre a competição que determina o melhor projeto
da \ac{ages}. O julgamento é baseado no que o projeto faz e em seu apelo de mercado,
e não na complexidade técnica da solução. Isso significa que um projeto visualmente
atraente e alinhado às tendências pode ter vantagem sobre um trabalho tecnicamente
mais elaborado e criterioso. Nosso projeto não era visualmente sofisticado, mas
seguia fielmente o que a stakeholder solicitava. Isso não diminui em nada a
experiência incrível que foi participar dele. Apenas agradeço por ter feito parte de
um time tão dedicado e por tudo que aprendi ao longo dessa jornada.

\textbf{Lições aprendidas:}
\begin{itemize}
  \item \textbf{Positivo:} colaboração entre os diferentes níveis de \ac{ages}, apoio
    constante dos colegas mais experientes, entrega funcional ao final do semestre
    e crescimento técnico expressivo partindo do zero;
  \item \textbf{A melhorar:} a integração entre frontend e backend poderia ter sido
    iniciada mais cedo, o banco de dados deveria ter sido finalizado com mais
    antecedência, e o tempo dedicado ao aprendizado individual das tecnologias
    poderia ter sido mais estruturado desde o início.
\end{itemize}
